% !TEX root = ../algebra_02.tex

\section{Жорданова нормальная форма (формулировка теоремы)}

\begin{Definition}{Жорданова клетка и жорданова матрица}{}
	Пусть $k$ — поле, $\lambda \in k$, $n \in \mathbb{N}$. 
	\emph{Жордановой клеткой} порядка $n$ с собственным значением $\lambda$ называется квадратная матрица вида:
	\[
	J_n(\lambda) = \begin{pmatrix}
		\lambda & 0 & \dots & 0 \\
		1 & \lambda & \dots & 0 \\
		\vdots & \ddots & \ddots & \vdots \\
		0 & \dots & 1 & \lambda
	\end{pmatrix}
	\]
	\emph{Жордановой матрицей} называется блочно-диагональная матрица, блоки которой — жордановы клетки:
	\[
	J = \operatorname{diag}(J_{k_1}(\lambda_1), \dots, J_{k_m}(\lambda_m)).
	\]
\end{Definition}

\begin{Definition}{Жорданов базис}{}
	Пусть $A \in \operatorname{End} V$. Базис $E$ пространства $V$ называется \emph{жордановым базисом}, если матрица оператора в этом базисе $[A]_E$ является жордановой матрицей.
\end{Definition}

\begin{Theorem}{О жордановой нормальной форме}{}
	Пусть $A \in \operatorname{End} V$ и его характеристический многочлен $\chi_A$ раскладывается на линейные множители над полем $k$. Тогда:
	\begin{enumerate}
		\item В пространстве $V$ существует жорданов базис.
		\item Если $E$ и $E'$ — два жордановых базиса пространства $V$, то матрицы $[A]_E$ и $[A]_{E'}$ отличаются только порядком жордановых блоков.
	\end{enumerate}
\end{Theorem}

\begin{Remark}{Связь с характеристическим многочленом}{}
	Если в жордановом базисе $E$ матрица оператора имеет вид $[A]_E = \operatorname{diag}(J_{k_1}(\mu_1), \dots, J_{k_l}(\mu_l))$, то характеристический многочлен равен:
	\[
	\chi_A(x) = \pm (x - \mu_1)^{k_1} \dots (x - \mu_l)^{k_l},
	\]
	где собственные значения $\mu_1, \dots, \mu_l$ не обязательно различны (то есть одному корню может соответствовать несколько жордановых клеток).
\end{Remark}

\begin{Statement}{Матричная формулировка (Следствие)}{}
	Пусть матрица $A \in M_n(k)$ такова, что её характеристический многочлен $\chi_A$ раскладывается на линейные множители. Тогда:
	\begin{enumerate}
		\item Существует обратимая матрица $U \in GL_n(k)$ такая, что $U A U^{-1}$ — жорданова матрица.
		\item Две жордановы матрицы подобны тогда и только тогда, когда они отличаются только порядком клеток.
	\end{enumerate}
\end{Statement}